\chapter{RESEARCH METHODOLOGY AND SYSTEM DESIGN}
\label{chap:methodology}

This chapter presents the methodology and system design of the proposed
Vision-Language-Action (VLA) framework for robot manipulation under constrained
resources. The chapter focuses on three project priorities: (1) constructing an
end-to-end VLA pipeline that is practical for commercial setups, (2) satisfying data,
hardware, and safety constraints for Dobot CR5 deployment, and (3) defining a clear
evaluation protocol across simulation and real-world domains.

To improve traceability between design and evidence, the chapter is structured from
requirements to implementation and then to evaluation. Each design decision is
explicitly linked to one or more project constraints introduced in
Chapter~\ref{chap:introduction}.


%─────────────────────────────────────────────────────────────────────────────
\section{METHODOLOGY OVERVIEW}
%─────────────────────────────────────────────────────────────────────────────

The research workflow is organized as an implementation pipeline from communication
infrastructure to model refinement:

\begin{enumerate}[label=\arabic*)]
    \item \textbf{AI-to-Robot Communication Interface}: build a module that converts
    Dobot TCP/IP command primitives into a machine-learning-friendly composable action
    format and vice versa.

    \item \textbf{Dataset Collection}: construct training data from two sources:
    (a) simulation trajectories in Isaac Lab exported to LeRobot format, and
    (b) real-world human demonstrations on the physical robot.

    \item \textbf{Model Development}: implement and integrate multimodal perception,
    language grounding, and action prediction modules into one end-to-end model.

    \item \textbf{Training Objective Design}: define losses and optimization settings
    for stable learning under low-data and limited-memory conditions.

    \item \textbf{Inference and Model Analysis}: inspect prediction behavior,
    latency, failure modes, and sim-to-real transfer quality.

    \item \textbf{Model Improvement}: apply targeted refinement based on analysis,
    including data balancing, prompt/interface tuning, and architecture adjustments.
\end{enumerate}

At a high level, the learning problem is to estimate a policy
\(\pi_\theta\) that maps multimodal context to controller-executable actions:
\[
a_t = \pi_\theta(I_t, \mathcal{U}, h_t),
\]
where \(I_t\) is the RGB observation, \(\mathcal{U}\) is the natural-language
instruction, and \(h_t\) is optional history/state context. The methodology seeks to
maximize task success under three constraints: limited demonstrations, limited memory,
and safety-constrained execution.


%─────────────────────────────────────────────────────────────────────────────
\section{PROBLEM-DRIVEN DESIGN REQUIREMENTS}
%─────────────────────────────────────────────────────────────────────────────

The system design is directly derived from the project constraints. Unlike large-scale
VLA studies that rely on broad datasets and high-end servers, this thesis targets a
minimal commercial setup. Therefore, design decisions are driven by three major
constraints.

\subsection{Dataset Quantity and Variety Constraint}

The available dataset is intentionally limited: approximately one hour of total robot
motion with only three task families. This is significantly smaller than many VLA
pipelines that use more than ten hours of demonstrations and broader task diversity.
As a result, the methodology prioritizes sample efficiency, high data quality,
cross-domain consistency (sim and real), and curriculum-style task coverage.

In practical terms, this means the pipeline avoids data-hungry assumptions and instead
prioritizes representation reuse, careful trajectory normalization, and targeted
coverage of failure-prone scenarios.

\subsection{Hardware Constraint (Commercial GPU)}

The model architecture must fit within practical hardware specifications,
specifically around 8 GB of VRAM. This restricts parameter count, sequence length,
and batch-size strategy, and motivates memory-aware training techniques such as mixed
precision, gradient accumulation, compact tokenization, and lightweight heads.

This constraint also affects experiment planning: ablations and evaluation schedules
must be designed to preserve comparability without requiring large parallel runs.

\subsection{Safety and Inference Standard Constraint}

For safe and repeatable operation on Dobot CR5, robot joint-position standards are
fixed before inference to define a consistent workspace and startup condition. The
exact joint preset is recorded in the experiment configuration file and should be
reported as
\texttt{[Joint preset values for Dobot CR5 to be specified in final version]}.
This standardization reduces run-to-run variance and lowers risk during autonomous
execution.

The safety layer is treated as a first-class system component rather than a post hoc
filter. This design choice is critical when model predictions are generated from
natural-language input, where ambiguity can propagate to unsafe control proposals.




%─────────────────────────────────────────────────────────────────────────────
\section{VISION-LANGUAGE-ACTION SYSTEM ARCHITECTURE}
%─────────────────────────────────────────────────────────────────────────────

\subsection{Input and Output Definition}

At each control cycle, the system receives:
\begin{itemize}
    \item an RGB observation $I_t$, and
    \item a natural-language instruction $\mathcal{U}$ from the user.
\end{itemize}

The system outputs a structured action command executable by the robot arm controller.
The output is designed to be composable with the TCP/IP communication interface so it
can be translated safely into Dobot command streams.

Operationally, the architecture separates semantic grounding from command realization.
This improves debuggability and allows independent validation of perception-language
alignment and action executability.

\subsection{VLM Core}

The VLM module performs multimodal understanding and instruction grounding. It
extracts intent from natural language, associates intent with current scene context,
and generates a high-level task state for downstream action production.

The VLM output is intentionally constrained to a structured intermediate task state
instead of free-form control text. This reduces ambiguity at the interface with the
action generator.

\subsection{Vision Expert}

The Vision Expert provides task-relevant visual features (object cues, target
relations, workspace context, and target hints) for robust action planning. It is
designed to improve perception reliability under object and scene variation.

To support low-resource training, the Vision Expert emphasizes transferable features
that remain stable across simulation and real camera conditions.

\subsection{Action Expert}

The Action Expert converts grounded task representations into executable robot
manipulation commands. The module is constrained by communication safety rules,
workspace limits, and robot kinematics, ensuring that generated actions remain valid
for runtime execution.

The Action Expert also acts as a control regularizer: it converts semantic intent into
bounded command trajectories that are compatible with controller timing and packet
format requirements.

\subsection{Module Coordination}

The three modules are connected through explicit intermediate representations:
\begin{enumerate}[label=\alph*)]
    \item \textbf{Language-grounded task state} from VLM,
    \item \textbf{Task-conditioned visual features} from Vision Expert,
    \item \textbf{Executable action sequence} from Action Expert.
\end{enumerate}

This separation improves maintainability and enables targeted optimization of each
module while keeping end-to-end behavior coherent.

In addition, the architecture introduces a protocol adapter between the Action Expert
and Dobot controller. This adapter maps model-level action semantics to TCP/IP command
payloads with explicit validation rules (range checks, rate checks, and command
structure checks) before dispatch.

The protocol adapter provides deterministic verification before hardware dispatch,
which is essential for reproducibility and for safe execution in real-robot tests.


%─────────────────────────────────────────────────────────────────────────────
\section{IMPLEMENTATION STRATEGY}
%─────────────────────────────────────────────────────────────────────────────

\subsection{Dataset Collection Strategy}

The data pipeline combines simulation and real-world trajectories:
\begin{enumerate}[label=\arabic*)]
    \item Generate simulation trajectories in Isaac Lab.
    \item Export trajectories to LeRobot-compatible format.
    \item Collect human demonstrations on the physical Dobot CR5.
    \item Normalize both sources into a unified schema of observation,
    instruction, and action tuples.
\end{enumerate}

This two-source design improves sample efficiency: simulation contributes scalable
coverage of trajectory patterns, while real demonstrations provide embodiment-specific
corrections for dynamics and sensing mismatch.

\subsection{Model Development and Training Strategy}

The system is trained and tuned in stages:
\begin{enumerate}[label=\arabic*)]
    \item Initialize the VLA backbone with pretrained multimodal knowledge.
    \item Adapt perception and action heads to robot-manipulation dynamics.
    \item Train with low-resource-aware objectives for stability and action validity.
    \item Perform end-to-end tuning to reduce module mismatch and command noise.
\end{enumerate}

The training objective combines command prediction and execution validity terms:
\[
\mathcal{L}_{\text{total}} = \lambda_1 \mathcal{L}_{\text{action}}
+ \lambda_2 \mathcal{L}_{\text{validity}}
+ \lambda_3 \mathcal{L}_{\text{alignment}},
\]
where \(\mathcal{L}_{\text{action}}\) supervises target actions,
\(\mathcal{L}_{\text{validity}}\) penalizes unsafe or non-executable outputs, and
\(\mathcal{L}_{\text{alignment}}\) preserves language-vision-task consistency.

\subsection{Inference and Model Analysis}

After training, inference behavior is analyzed in both simulation and physical setups,
including latency profile, command validity, failure taxonomy, and transfer behavior.
This analysis identifies whether errors are caused by perception ambiguity, language
misgrounding, or action decoding instability.

Failure analysis is organized by stage so corrective interventions can be applied with
minimal retraining overhead.

\subsection{Model Improvement Loop}

Based on analysis findings, iterative improvements are applied through data balancing,
task reweighting, communication-level safeguards, and architectural tuning.

This closed-loop process is central to low-resource development because each
improvement cycle must produce measurable gains without expensive data re-collection.

\subsection{Latency and Memory Optimization}

To satisfy real-time robot-control requirements, the implementation applies practical
optimization techniques, including mixed precision inference, efficient token/image
preprocessing, and lightweight module interfaces that remain feasible on 8 GB VRAM.

Together, these optimizations convert the architecture from a proof-of-concept model
into an executable pipeline for interactive robot operation.


%─────────────────────────────────────────────────────────────────────────────
\section{EVALUATION PROTOCOL FOR MANIPULATION TASKS}
%─────────────────────────────────────────────────────────────────────────────

\subsection{Evaluation Objectives}

Evaluation is designed to measure whether the system satisfies the project goals:
\begin{itemize}
    \item follows natural-language instructions,
    \item generalizes to unseen tasks,
    \item maintains stable manipulation output,
    \item satisfies real-time execution constraints.
\end{itemize}

The protocol deliberately separates robustness assessment (success-rate domain) from
transfer assessment (zero-shot domain), because these dimensions often improve at
different rates in constrained-data training.

\subsection{Evaluation Domains}

The protocol consists of two main domains:
\begin{enumerate}[label=\arabic*)]
    \item \textbf{Success-Rate Evaluation}:
    assess task completion under standard conditions.
    \item \textbf{Zero-Shot Evaluation}:
    assess transfer to unseen combinations without additional fine-tuning.
\end{enumerate}

Each domain is evaluated in both environments:
\begin{itemize}
    \item On simulation,
    \item On real robot.
\end{itemize}

This matrix of \{success-rate, zero-shot\} \(\times\) \{sim, real\} is used to
distinguish pure policy quality from embodiment transfer effects.

\subsection{Task Categories}

The benchmark includes representative robot manipulation tasks such as:
\begin{enumerate}[label=\arabic*)]
    \item pick-and-place,
    \item stacking,
    \item drawer or container interaction,
    \item multi-step object sequencing.
\end{enumerate}

\subsection{Evaluation Metrics}

The following metrics are used:
\begin{itemize}
    \item \textbf{Task success rate}: whether generated actions satisfy the command
    objective in each task.
    \item \textbf{Zero-shot success rate}: performance on unseen combinations relative
    to standard seen conditions.
    \item \textbf{Action validity}: percentage of outputs that are executable by the
    robot controller.
    \item \textbf{Latency}: end-to-end action generation time per command cycle.
\end{itemize}

For comparative reporting, unseen-task retention is additionally summarized as:
\[
\text{Generalization Ratio} =
\frac{\text{Unseen Task Success}}{\text{Seen Task Success}} \times 100.
\]


%─────────────────────────────────────────────────────────────────────────────
\section{CHAPTER SUMMARY}
%─────────────────────────────────────────────────────────────────────────────

This chapter defined a practical VLA methodology aligned with low-resource
implementation goals. The proposed system integrates multimodal reasoning and robot
action generation through a standardized AI-to-robot interface. The design explicitly
addresses limited dataset scale, 8 GB VRAM hardware constraints, and Dobot CR5 safety
standards, and establishes an evaluation protocol across simulation and real-world
domains for both standard success-rate and zero-shot tests.

With this methodology in place, Chapter~\ref{chap:experiments} evaluates whether the
designed pipeline delivers measurable gains in action validity, instruction following,
and transfer performance under the targeted resource constraints.
